\chapter{Tổng quan}
\label{chap:overview}

\section{Kiến thức Cơ sở về Hệ thống Dẫn động Cơ khí}
\label{sec:overview-mech}

\subsection{Hệ thống Dẫn động Thùng Trộn}

Đề bài cơ kỹ thuật (Phương án~2) yêu cầu thiết kế hệ thống dẫn động thùng trộn
với sơ đồ động học như sau:

\[
  \text{Động cơ điện} \xrightarrow{\text{nối trục đàn hồi}}
  \text{HGT côn--trụ 2 cấp} \xrightarrow{\text{xích ống con lăn}}
  \text{Thùng trộn}
\]

Các thông số thiết kế được cho sẵn theo Phương án~2:

\begin{table}[H]
  \centering
  \caption{Thông số thiết kế đầu vào}
  \label{tab:input-params}
  \begin{tabular}{llcc}
    \toprule
    \textbf{Ký hiệu} & \textbf{Thông số}                    & \textbf{Giá trị} & \textbf{Đơn vị} \\
    \midrule
    $P$              & Công suất trên thùng trộn            & 7,5              & kW              \\
    $n$              & Số vòng quay trên trục thùng trộn    & 75               & v/p             \\
    $L$              & Thời gian phục vụ                    & 10               & năm             \\
    --               & Chế độ vận hành                      & \multicolumn{2}{l}{Quay 1 chiều, tải va đập nhẹ} \\
    --               & Lịch làm việc                        & \multicolumn{2}{l}{300 ngày/năm, 2 ca/ngày, 8 h/ca} \\
    \bottomrule
  \end{tabular}
\end{table}

Hệ thống bao gồm các bộ phận truyền động sau, mỗi bộ phận sẽ được hỗ trợ bởi một module tính toán trong ứng dụng:


\begin{itemize}
    \item \textbf{Động cơ điện ba pha không đồng bộ}: nguồn động lực cho toàn hệ.
    \item \textbf{Nối trục đàn hồi}: kết nối trục động cơ với trục vào hộp giảm tốc.
    \item \textbf{Hộp giảm tốc côn--trụ hai cấp}: cấp nhanh dùng bánh răng côn,
            cấp chậm dùng bánh răng trụ răng thẳng.
    \item \textbf{Bộ truyền xích ống con lăn}: nối trục ra hộp giảm tốc với thùng trộn.
\end{itemize}


\subsection{Các Module Tính toán Cần Hỗ trợ}

Quy trình tính chọn động cơ gồm các bước \cite{LeNhatTruong2026}

\subsubsection*{Module 1 -- Tính Chọn Động cơ và Phân phối Tỉ số Truyền}
\begin{itemize}
    \item Tính chọn động cơ điện
    \item Phân phối tỷ số truyền
\end{itemize}

\subsubsection*{Module 2 -- Tính bộ truyền xích}
\begin{itemize}
    \item Tính thông số thiết kế bộ truyền xích
    \item Kiểm nghiệm bộ truyền
\end{itemize}

\subsubsection*{Module 3 -- Tính bộ truyền bánh răng}
\begin{itemize}
    \item Tính bộ truyền kín cấp nhanh
    \item Tính bộ truyền kín cấp chậm
\end{itemize}


\section{Các Công nghệ Sử dụng}
\label{sec:overview-tech}

\subsection{Nền tảng Phát triển Ứng dụng Di động}

Ứng dụng được xây dựng trên nền tảng \textbf{React Native} phiên bản 0.81 kết hợp với \textbf{Expo SDK 54}.
React Native cho phép viết một codebase duy nhất bằng TypeScript và biên dịch ra ứng dụng gốc (native)
cho cả Android và iOS, đáp ứng yêu cầu đa nền tảng của đề tài.

\textbf{Expo Router v6} được sử dụng để quản lý điều hướng theo cơ chế \textit{file-based routing}:
mỗi tệp trong thư mục \texttt{src/app/} tương ứng với một màn hình,
tương tự cách Next.js hoạt động trên web.
Cơ chế này giúp giảm boilerplate và dễ dàng mở rộng thêm màn hình mới.

\textbf{NativeWind v4} (Tailwind CSS cho React Native) được dùng để định kiểu giao diện
thông qua các utility class, đạt giao diện nhất quán mà không cần viết \texttt{StyleSheet} thủ công.

\subsection{Cơ sở Dữ liệu Cục bộ}

Dữ liệu cục bộ được lưu trữ trong \textbf{SQLite} thông qua thư viện \texttt{expo-sqlite}.
ORM \textbf{Drizzle v0.44} được dùng để định nghĩa schema và thực hiện truy vấn kiểu an toàn
(type-safe), tránh lỗi SQL thủ công.

Dữ liệu bảng tra cứu động cơ (catalog ĐK và 4A) được đóng gói sẵn vào ứng dụng
dưới dạng seed data TypeScript.
Khi ứng dụng khởi động lần đầu, hệ thống tự động chạy migration tạo bảng
và nạp dữ liệu động cơ vào SQLite, đảm bảo tính năng tra cứu hoạt động hoàn toàn offline.

\subsection{Quản lý Trạng thái và Giao tiếp Liên Module}

Trạng thái liên module được quản lý bằng đối tượng \textbf{globalNavigationState} --
một plain object toàn cục được chia sẻ giữa các màn hình tính toán.
Khi Module~1 hoàn thành, kết quả được ghi vào \texttt{globalNavigationState.gearResult}
và \texttt{chainResult}; Module~2 và 3 đọc lại các giá trị này khi màn hình được điều hướng đến.

Dữ liệu phiên tính toán và tài khoản người dùng được đồng bộ với backend
qua thư viện \textbf{Axios}.
JWT access token và refresh token được lưu cục bộ bằng \textbf{AsyncStorage},
cho phép tự động khôi phục phiên đăng nhập.

\subsection{Backend}

Backend được xây dựng bằng \textbf{Java Spring Boot} với cơ sở dữ liệu \textbf{PostgreSQL},
chịu trách nhiệm xác thực người dùng và lưu trữ phiên tính toán trên server.
Phần backend nằm ngoài phạm vi của báo cáo này; client mobile giao tiếp với backend
thông qua REST API được cấu hình qua biến môi trường \texttt{EXPO\_PUBLIC\_API\_URL}.
